% Backup of body.tex
% Created by Copilot assistant before synonym/paraphrase edits.
% Original content preserved.

% !Mode:: "TeX:UTF-8"

\iffalse
\bibliography{static/reference/reference.bib} % 欺骗latextools获取bib文件
\fi

%? ************************************************************* %?
%?                                                               %?
%?                   您将在此处完成正文内容书写                      %?
%?                                                               %?
%? ************************************************************* %?


% 全局body字体
\xfbody

\chapter{引言}
        \section{研究背景和意义}
        \xfbody{根据商务部发布的《中国新电商发展报告（2025）》\cite{spb2024_bulletin_20250522}，2024年我国网络零售额达到15.5万亿元，同比增长7.2\%，对社会消费品零售总额贡献约1.7个百分点，我国已连续12年保持全球最大网络零售市场的地位，相关数据见图\ref{fig:kuaidi1}。电商规模的持续扩大使仓储与物流环节逐步演变为整体效率提升的主要瓶颈。依靠人工完成的仓储作业在效率和响应速度上已无法匹配不断攀升的订单量，这一现实促使学术界和产业界将目光投向智能仓储中的自主移动机器人。

\begin{figure}[H]
        \centering
        \includegraphics[width=0.5\textwidth]{kuaidi1.png}
        \caption{2020-2024年人均用邮支出、快递支出和快递使用量情况}
        \label{fig:kuaidi1}
        \end{figure}

        在智能仓储领域，常见的智能仓储移动机器人有如下几类：自动导引车（AGV）、自主移动机器人（AMR）、无人搬运车（UGV）和智能搬运机器人等。AGV是最早实现商业化应用的智能仓储移动机器人，主要依靠磁条、二维码等环境标识进行导航，适用于路径相对固定、环境较为简单的场景；AMR则具备更强的环境感知和自主决策能力，能够在动态变化的环境中灵活导航，适合复杂多变的仓储环境；UGV通常指具备越野能力的无人搬运车，适用于户外或地面不平整的场景；智能搬运机器人则强调在搬运过程中对货物的智能识别和处理能力，能够实现从搬运整组货架到单件商品的操作升级。

        自动导引车（Automated Guided Vehicle, AGV）是自主移动机器人的典型代表，它集自主导航、路径规划、货物搬运与障碍规避于一体，可在复杂仓储场景中独立完成物流作业。图\ref{fig:AGV}展示了AGV小车的工作状态，在仓储系统里AGV主要负责运输和调度。相较于纯人工搬运方式，AGV能够明显提高作业效率与准确率，同时削减人力开支并改善安全状况。从市场数据来看，2019年到2024年中国AGV行业规模从48.27亿元攀升至95.49亿元，足以说明市场对智能搬运装备的旺盛需求。

\begin{figure}[H]
        \centering
        \includegraphics[width=0.5\textwidth]{AGV.png}
        \caption{AGV小车工作示意图}
        \label{fig:AGV}
        \end{figure}
        }

        \section{国内外智能仓储移动机器人研究综述}
        \xfbody{
                
                \subsection{智能仓储移动机器人国外研究现状}
                \xfbody{
                国外在仓储机器人领域的研发和商业化比起国内落地起步更早，成熟度也更高。亚马逊的Kiva系统是最具代表性的案例：2012年亚马逊斥资7.75亿美元收购Kiva公司，此后在全球仓库中批量部署Kiva机器人，以货架到人模式取代了传统人到货架的拣选流程，拣选效率提高了2至3倍，这一实践至今仍是AGV商业化运营的标杆\cite{李成进2016智能移动机器人导航控制技术综述}。欧洲方面，德国Magazino公司深耕精细化仓储操作，其TORU机器人不仅能够搬运货架，还可以借助视觉识别从货架上直接抓取单件商品，完成了从搬运整组货架到拾取单品的技术跨越。美国波士顿动力公司长期聚焦仿生机器人方向，Atlas人形机器人和Spot四足机器人在复杂地形行走、动态平衡控制等领域展现出较高水准。

                在算法维度，国外的研究范围较为宽泛。路径规划领域，Hybrid A*算法由于能够生成满足机器人运动约束的可行路径，在实际产品中得到了广泛采纳\cite{nemec2019improving}；多机调度领域，多智能体路径规划（MAPF）算法为多台AGV的协同运作提供了理论支撑。最近两年的研究焦点转向深度强化学习的落地应用，借助DQN、PPO等算法使机器人在与环境的交互过程中自主习得导航策略，从而应对仓储场景的动态变化。

                梳理国外移动机器人的演进脉络，不难发现一条从理论突破到工程实践再到场景拓展的清晰路径：1961年第一台工业机器人Unimate问世，1966年自主移动机器人Shakey诞生，2012年Kiva实现大规模商业化部署，近年Atlas在动态环境中展现出令人瞩目的性能——每一步跨越都离不开理论探索与工程落地的相互驱动。

}
                }
                \subsection{智能仓储移动机器人国内研究现状}
                \xfbody{国内对AGV的研究可追溯到20世纪70年代末期，彼时的工作重心放在引进和仿制国外设备上，自主技术积淀不多。21世纪以来，制造业转型升级叠加电商行业的高速增长，使AGV逐步进入产业界和学术界的视野。2012年亚马逊收购Kiva Systems一事尤其具有标志性意义，国内企业由此认识到仓储机器人的巨大发展空间，京东、阿里、海康威视等公司先后布局这一赛道\cite{谢轲晗2025数字赋能制造企业与物流企业融合创新发展的模式及路径研究}，在导航、避障、路径规划和货物识别等核心环节展开攻关，从传感器、控制器到整机系统的国产化产业链也随之逐步成形。

                在应用端，国内AGV已由最初的汽车制造和烟草仓储延伸至电商物流、医药配送、电子装配等诸多行业\cite{田璐2026ai}。导航技术的迭代同样值得关注：早期产品多依赖磁条或二维码引导，需要在地面铺设标识，部署费用偏高且灵活性不足；近几年激光雷达和视觉传感器价格持续走低，基于激光SLAM和视觉SLAM的自然导航方案逐步占据主流地位，这类方案无需改造现有环境即可部署，其中视觉SLAM因传感器成本更低而尤其受到产业界青睐\cite{沈博闻2014仓储物流机器人集群的智能调度和路径规划}。在算法层面，国内研究者围绕路径规划和多机调度先后提出了改进A*、遗传算法、蚁群算法、粒子群优化等多种策略，积累了一定的理论成果。

                坦率地讲，国内AGV在核心算法的成熟度、系统可靠性以及大规模部署经验上与国际先进水平之间还存在差距，不过这一差距正在逐步收窄。

}
        \section{国内外智能仓储移动机器人路径规划算法研究现状}
        \xfbody{{

        在国外研究方面，M Rybczak等人梳理了移动机器人控制中机器学习算法面临的计算资源、实时决策、环境适应性和安全可靠性等挑战\cite{rybczak2024survey}。MFR Lee等人将深度Q网络（DQN）和双深度Q网络（DDQN）引入移动机器人导航与避撞学习\cite{lee2022mobile}。Y Dai等人剖析了传统鲸鱼优化算法及其改进方法在移动机器人路径规划中的不足\cite{dai2023novel}。S Das等人提出了自适应随机梯度下降线性回归（ASGDLR）算法，用于识别自主移动机器人左右转向运动特征\cite{das2022machine}。B Tao等人提出了具有随机扰动策略的双种群粒子群优化算法，以改善传统粒子群优化在约束问题中的局限性\cite{tao2024mobile}。T Guo等人提出了一种基于改进A*和动态窗口法（DWA）融合的AGV路径规划算法，旨在解决传统A*搜索效率低、不考虑AGV尺寸和转弯次数过多，以及DWA易陷入局部最优的问题\cite{guo2023automated}。K Fransen等人提出了一种改进的A*算法启发式方法，可用于在具有转向成本的几何图中找到最低成本路径\cite{fransen2023efficient}。T Heinemann等人提出了一种路径规划概念及评估方法，用于改进向量场直方图（vector field histogram）和动态窗口（dynamic window）算法，以实现平滑轨迹生成\cite{heinemann2019generating}。

        在国内研究方面，学者们围绕仓储机器人路径规划与轨迹跟踪进行了大量探索。黄晓宇等人设计了一种模型预测控制（MPC）与微分先行比例-积分-微分（PID）协同的双闭环控制策略，用于麦克纳姆轮全向移动平台的轨迹跟踪控制\cite{黄晓宇2023基于}。曹梦龙等人梳理了机器人路径规划的基本步骤，并概述了群体智能优化算法在路径规划中的应用\cite{cao2023robot}。白宇鑫等人提出了改进的哈里斯鹰优化算法以解决机器人路径规划问题，并剖析了HHO算法在收敛性和局部最优方面的局限性\cite{白宇鑫2025基于改进哈里斯鹰算法的机器人路径规划研究}。赖荣燊等人探讨了传统A*算法与动态窗口法的优缺点及其改进方向\cite{赖荣燊2024融合改进}。张瑞等人提出了RRT*与DWA融合的路径规划方法，旨在提升复杂动态环境下移动机器人路径规划的安全性和最优性\cite{张瑞2024融合}。肖金壮等人设计了一种面向室内自动导引车（AGV）路径规划的改进蚁群算法，旨在解决传统算法在大规模和复杂环境中全局搜索效率差、收敛速度慢以及路径转弯次数过多和不够平滑的问题\cite{肖金壮2022一种面向室内}。张中伟等人提出了一种基于动态优先级策略的多AGV无冲突路径规划方法，用于柔性制造车间的高效平稳生产。该方法通过动态调整AGV的优先级，优化路径规划过程，减少冲突和等待时间，提高整体生产效率\cite{张中伟2021基于动态优先级策略的多}。张天瑞等人探讨了针对基本蚁群算法在机器人路径规划过程中存在的路径转弯角度大、易陷入局部极小值和收敛速度慢等问题所进行的改进策略\cite{张天瑞2022面向机器人全局路径规划的改进蚁群算法研究}。李玉清等人设计了一种针对自动导引车（AGV）路径规划的改进动态窗口法（DWA）算法，该算法克服了原始DWA存在的振荡现象，并融合了人工势场法（APF）以提升路径规划的安全性和效率\cite{李玉清2025一种动态窗口法和人工势场法融合的}。
        
        综合来看，国内外在智能仓储移动机器人路径规划算法方面已经形成了较为系统的研究基础，研究脉络主要经历了从传统启发式搜索到融合运动学约束的规划方法，再到结合机器学习与多机协同的智能化发展过程，后续研究重点将继续聚焦高效规划、动态避障与复杂场景下的协同决策。

        }
        \section{研究方法与手段}
                \xfbody{AGV的本体形态日趋多元化，旨在适配不同的仓储场景和搬运需求。本文所研究的对象是一款具备举升功能的自主搬运机器人，它集载货、升降与运输于一体，依靠举升机构完成货物装卸作业，可在一般工业和仓储环境中使用，且具有较好的扩展性和场景适应能力。

                底盘设计环节着重考量工作方式、电机选型与平台布局等要素。整套系统包含升降机构、行走装置和避障模块三个部分：升降机构负责载物台的抬升与下降，行走装置驱动车辆移动，避障模块借助视觉或非视觉传感器探测前方障碍物，从而增强运行过程中的安全性。升降机构选用叉剪式结构，具备传动平稳、结构紧凑的优点；若面对大载荷工况，也可改用桅柱式或套筒式等替代方案。本文对差速驱动底盘进行运动学建模，差速底盘一般由两侧独立驱动轮和若干万向轮组成，依靠调节左右轮的速度差来实现行驶和转向。控制策略选用PID控制器，利用比例、积分与微分三种作用来缩减跟踪误差，图\ref{fig:pid}呈现了系统的控制框架。

                \begin{figure}[H]
                \centering
                \includegraphics[width=0.7\textwidth]{pid.png}
                \caption{PID控制示意图}
                \label{fig:pid}
                \end{figure}


本文的主要工作包括：

1. 智能仓储移动机器人总体方案设计：提出适用于智能仓储的智能仓储移动机器人系统总体方案，确定差速底盘、路径规划与避障策略。

2. 机械结构设计：完成载物与升降机构、行走与避障模块的结构设计，并开展强度校核与有限元分析。

3. 电气系统设计：完成电机、传感器与控制电路的选型与原理图设计。

4. 运动学建模与仿真：建立差速底盘模型，并利用MATLAB/Simulink进行轨迹仿真验证。

5. 运动控制研究：设计并验证基于PID的双轮差速控制策略，实现基本运动控制性能验证。

6. 结论与展望：总结研究成果，指出不足并提出后续改进方向。

}
\chapter{智能仓储移动机器人总体设计方案}
    \section{仓储移动机器人选型}
        \xfbody{智能仓储移动机器人根据功能定位和应用场景的不同可以分为多种类型。常见机器人类型的核心特点与应用场景对比见表\ref{tab:robot_types}。

\begin{table}[H]
        \wuhao
        \centering
        \caption{智能仓储移动机器人类型对比}
        \label{tab:robot_types}
        \begin{tabular}{p{2cm}p{4cm}p{4cm}}
        \toprule
        机器人类型 & 核心特点 & 典型应用场景 \\
        \midrule
        自主移动机器人(AMR) & 具备自主导航与动态避障能力，无需改造环境，灵活性最高。 & 人机混场的动态环境、复杂搬运、高频拣选等。 \\
        自动引导车(AGV) & 需沿预设固定路径行驶，技术成熟，成本相对较低。 & 固定路线、大规模、重复性的物料搬运（如Kiva类机器人）。 \\
        仓储无人机(UAV) & 在空中飞行，快速抵达地面机器人难以企及的高处。 & 高位货架库存盘点、轻量级紧急物品搬运等。 \\
        \bottomrule
        \end{tabular}
        \end{table}

现代仓储环境对货物运输效率和灵活性提出了较高的要求。AGV的运行逻辑相对固定，适用于路径相对简单、环境较为稳定的场景。因其研究较早，技术成熟，在业务流程固定、追求高性价比的大规模搬运场景中，AGV依然是主流选择。本文选择研究的智能仓储移动机器人属于AGV范畴，具备举升功能，能够在货架密集的环境中自主行驶，执行物料的自动搬运与装卸任务。

        }
        \section{设计需求和技术指标}
    \xfbody{电商物流和智能制造的快速推进，使依靠人工驾驶叉车或人力推车完成仓储搬运的传统方式在效率、安全性和人力开支等方面遭遇日益严峻的考验。为了满足智能仓储对空间利用率和作业灵活性的双重需求，本文设计了一款拥有自主导航能力的举升式移动机器人，它能够在货架密集的环境中自主行驶，执行物料的自动搬运与装卸任务。该系统以差速驱动\cite{刘梓博2025基于双变量限幅}底盘为基础，结合举升机构、避障传感器和路径规划算法，实现从取货点到目标货架之间的安全自主行走。

该智能仓储移动机器人的主要技术指标如下：

自动载物功能：配备可升降的载物台，升降行程为0.2m，可满足不同货架高度的搬运需求。

尺寸指标：外形设计需保证小车在狭窄通道与密集货架之间灵活通行。

结构组成：整车包括载物装置、行走装置、避障装置等核心模块，通过多种传感器增强环境感知能力。

载重能力：最大承载重量不超过25kg，适用于轻型物料的搬运场景。

智能避障与导航：通过超声波传感器实现中远距离障碍物检测，红外传感器实现地面循迹线检测，保证在复杂环境下的安全运行。

系统优势：空间利用率高、灵活性强、自动化程度高，能提升仓储作业效率，降低人工成本。

}


        \section{驱动选型}
        \xfbody{AGV的驱动方式受最大载荷、轮系数量、转向方式和驱动形式等多重因素制约，这些因素相互耦合，直接左右车辆的运动性能和适用范围，因而在设计阶段必须从整体系统视角出发，统筹兼顾结构、控制、成本和应用需求，才能得出合理的驱动方案。就目前而言，AGV的驱动形式大致可归为三类：单轮驱动、差速驱动和全方位驱动。

（1）单轮驱动

单轮驱动的特点在于一个驱动轮同时兼顾驱动和转向两项功能，从动轮对称布置在车体轴线两侧。驱动轮通常为舵轮，从动轮多为万向轮，如图\ref{fig:单轮驱动}所示。这种结构简洁、可靠性较高、维护费用低，适用于载荷不大、路径相对固定的场合。它能够实现前进、后退和左右转向（转角不超过$90^\circ$）。不过由于转弯半径偏大，灵活性欠佳，在空间局促或需要频繁转向的环境中表现不理想。

\begin{figure}[H]
        \centering
        \includegraphics[width=0.5\textwidth]{danlunqudong.png}
        \caption{单轮驱动示意图}
        \label{fig:单轮驱动}
        \end{figure}
        （2）差速驱动

差速驱动由两个驱动轮搭配两个或四个从动万向轮构成，左右驱动轮对称分布在车体轴线两侧。转向功能依靠改变两侧驱动轮的速度差来达成，如图\ref{fig:差速驱动}所示。这种驱动方式能够实现前进、后退和左右转向（转角超过$90^\circ$），转弯半径较小、灵活性出色，尤其适合频繁转向和狭窄通道的工况。不足之处在于驱动轮磨损较快，若是四轮结构还可能面临轮子悬空的问题，因此对地面平整度的要求较高，应用场景受到一定制约。

\begin{figure}[H]
        \centering
        \includegraphics[width=0.5\textwidth]{chasu.png}
        \caption{差速驱动示意图}
        \label{fig:差速驱动}
        \end{figure}
        （3）全方位驱动

全方位驱动的结构相对复杂，常见的实现方式有两种：一种是前后各布置一个舵轮并搭配两个从动万向轮，通过调节舵轮的角度和速度来完成转向，如图\ref{fig:全方位驱动}所示；另一种是采用四个麦克纳姆轮，分别独立控制四个轮子的转向和转速，借助速度矢量合成原理达成全方位运动。
        \begin{figure}[H]
        \centering
        \includegraphics[width=0.5\textwidth]{quanfangwei.png}
        \caption{全方位驱动示意图}
        \label{fig:全方位驱动}
        \end{figure}
        这种驱动方式可以在车身姿态保持不变的前提下朝任意方向行驶，灵活性极高，适合对空间利用率要求苛刻的场景。然而其结构较为复杂，对电机性能和控制精度的要求很高，开发难度和制造成本也随之攀升。

（4）设计选择与应用场景

结合功能需求，本次设计的AGV小车外形选用椭圆形，这种造型既能保证体积小巧美观，又能在左转或右转时尽量维持相同的行驶空间，特别适用于$90^\circ$直角转弯工况；椭圆形轮廓有助于减小转弯包络面积、提升狭窄通道通过性。需要说明的是，转弯半径主要由差速轮系和控制策略决定，而非仅由外形决定。在综合权衡结构复杂度、成本和可靠性之后，本次设计最终选定差速驱动作为举升式AGV小车的运行动力方案，原因在于该方案在确保灵活性的同时兼顾了稳定性和经济性。

}    
\chapter{智能仓储移动机器人机械结构设计}
    \section{核心部件结构设计}
    \xfbody{本节涉及移动底盘、驱动轮、升降机构和罩壳四个部分。第二章已就AGV小车的总体结构和驱动选型做了阐述，下面将从结构布局、材料选用、工艺路线和功能实现等方面，对各核心部件的结构设计展开详细探讨。

\subsection{移动底盘}
        \xfbody{移动底盘采用差速驱动方式，核心部件包括两个独立控制的驱动轮和四个辅助万向轮。差速驱动轮分置于底盘两侧，各自由独立电机驱动，依靠调节左右轮速度差来实现转向和行驶控制。辅助万向轮分布在底盘四角，提供额外的支撑和稳定性，保障智能仓储移动机器人在复杂仓储环境中的平稳运行。底盘结构采用爚接鈢框方案，兼具良好的强度和稳定性，同时预留了充足的空间用于安装电气系统和传感器。如图\ref{fig:dipan}所示，                    \begin{figure}[H]
        \centering
        \includegraphics[width=0.5\textwidth]{dipan.png}
        \caption{差速驱动底盘示意图}
        \label{fig:dipan}
        \end{figure}
        }
        \subsection{驱动轮}     
        \xfbody{驱动轮是AGV小车的关键部件之一，如图\ref{fig:qudonglun}所示。在该装置中，驱动部分的壳体选用45号钢板制造，厚度为2mm。45钢兼具较好的强度和韧性，既能保证整体刚性，又具备一定的减振能力。壳体采用弯曲与焊接工艺成型，在保证结构整体性的同时也控制了制造成本。驱动级的两个电机分别固定在壳体的左右两侧，形成对称布局，有助于车辆在运行中保持平衡。电机的另一端与驱动装置底座相连并固定，确保动力传递的可靠性。

\begin{figure}[H]
        \centering
        \includegraphics[width=0.5\textwidth]{qudonglun.png}
        \caption{驱动轮示意图}
        \label{fig:qudonglun}
        \end{figure}

驱动装置的底座处于整个结构的核心位置，其顶部装有直线轴承模块。直线轴承的作用在于保证支撑杆在上下运动时保持同心度，防止因偏移引发的振动或磨损。支撑杆安装在直线轴承模板上，构成稳定的导向结构，使整个驱动系统在工作过程中保持平稳。

在减振结构方面，装置设置了预紧螺母、支撑立轴、车身骨架连接板、固定板、阻尼弹簧和直线轴承板等关键部件。支撑立轴安装在直线轴承的中间位置以确保同心度，其顶部装有预紧螺母，用于减震器的固定和调节；下端通过圆柱销与驱动装置底座连接，形成可靠的约束。阻尼弹簧被压缩在车架连接板与固定板之间，能够在车辆运行中吸收冲击和振动，进而增强整体的减震效果。

该驱动装置在保持结构简洁的同时兼顾了减震性能与稳定性。

}
        \subsection{旋转机构部分}
        \xfbody{旋转机构是智能仓储移动机器人的重要功能模块之一，主要承担载物台的旋转任务，以便适应不同方向的搬运需求。如图\ref{fig:shengjiang}所示。

\begin{figure}[H]
        \centering
        \includegraphics[width=0.5\textwidth]{旋转机构.png}
        \caption{旋转机构示意图}
        \label{fig:shengjiang}
        \end{figure}
        该旋转机构采用电机驱动，经由齿轮传动带动载物台旋转。电机安装在底盘内部，通过齿轮连接，确保动力传递的稳定性和效率。旋转机构的设计充分考虑了载物台的重量和旋转惯量，选用高扭矩电机以保证满载工况下的平稳旋转。机构还配备了限位开关和编码器，用以实现旋转角度的精确控制和反馈，保障智能仓储移动机器人在执行搬运任务时能够准确调整载物台方向，从而提升作业效率和安全性。

}
        \subsection{罩壳设计}
        \xfbody{罩壳设计是智能仓储移动机器人不可或缺的组成部分，其主要职能是保护内部机械结构和电气系统免受外部环境干扰，同时兼顾美观和操作便利。如图\ref{fig:zhaoke}所示。

\begin{figure}[H]
        \centering
        \includegraphics[width=0.5\textwidth]{zhaoke.png}
        \caption{罩壳设计示意图}
        \label{fig:zhaoke}
        \end{figure}
        罩壳采用轻质高强度材料制成，具备良好的耐冲击性能和防护能力，可有效保护智能仓储移动机器人的内部组件免受碰撞、灰尘和湿气侵蚀。罩壳设计还兼顾了操作便利性，配有操作面板和指示灯，便于操作人员进行控制和监控。罩壳外形简洁美观，符合现代工业设计的审美标准，同时也能提高智能仓储移动机器人在仓储环境中的辨识度和安全性。                    }
        }
    \section{总体结构设计}
    \xfbody{
    在前述核心部件设计的基础上，对举升式智能仓储移动机器人的整体机械结构进行概述。智能仓储移动机器人总体机械结构的三维模型如图\ref{fig:agv_zongti}所示。

\begin{figure}[H]
        \centering
        \includegraphics[width=0.55\textwidth]{agv总体机械.png}
        \caption{智能仓储移动机器人总体机械结构三维模型}
        \label{fig:agv_zongti}
        \end{figure}

如图所示，智能仓储移动机器人总体结构按照功能分为四个层次：底盘承载层、驱动行走层、举升机构层和载物平台层。底盘由厚度15\,mm的Q235A碳素结构钢板焊接而成，外形呈椭圆形（长轴约1000\,mm，短轴约700\,mm），板面上预设了驱动轮、万向轮、电气仓和传感器等安装孔位。驱动行走层包含两个独立驱动轮和四个万向轮，驱动轮对称分布在底盘短轴两侧，由SMC80S伺服电机驱动；四个万向轮装在底盘四角，起辅助支撑和自由转向作用。举升机构层位于底盘上表面正中，行程200\,mm，升降期间载物台保持水平。载物平台层位于整车最上方，底部集成旋转机构，可绕竖直轴旋转$\pm90^\circ$。

整体外形尺寸满足设计限值要求，额定载重25\,kg。重心位置靠近驱动轮轴线，有利于行驶稳定性。各层之间依靠螺栓和定位销进行刚性装配，层间通过标准化接口相互配合，各部件的详细结构设计见第三章。
    }
    \section{结构强度校核和有限元分析}
        \subsection{力学分析选择}
        \xfbody{依据上述设计和计算结果，选定智能仓储移动机器人浮动支撑梁作为受力分析对象。智能仓储移动机器人本体重量为200 kg，额定载荷为25 kg，故满载情况下系统总重量为：

\begin{equation}
        W = (200 + 25) \cdot 9.81 \approx 2205 \, \text{N}
        \end{equation}

AGV轮系为“2个驱动轮+4个万向轮”。在结构初步校核中，为便于计算并取偏保守估计，将主要垂向载荷按两侧主承载梁分担处理，万向轮仅作为辅助支撑与导向。假设两侧主承载梁受力分布均匀，则单根支撑梁所承受的力可按总重量的一半估算，即：

\begin{equation}
        F = \frac{2205}{2} \approx 1102.5 \, \text{N}
        \end{equation}

在实际设计中，为保证安全性与可靠性，通常需要考虑安全系数。本文此处属于静载工况下的方案级初步校核，两侧主承载梁按均分假设各承担总载荷的一半，同时取1.02的名义放大系数以覆盖重心偏移等小量不确定性。后续工程定型阶段应依据GB/T 3811等规范重新核算安全系数。单根支撑梁的设计受力可近似为：

\begin{equation} 
        F_{\text{max}} \approx 1102.5 \times 1.02 \approx 1125 \, \text{N}
        \end{equation}

因此，在后续的结构设计与强度校核过程中，将以 1125 N 作为单根支撑梁的最大受力值进行分析。这一数值将作为梁的截面选择、材料强度校核以及疲劳寿命分析的重要依据。

}
        \subsection{支撑材料选择}
        \xfbody{在材料选择方面，综合权衡了结构强度、安全系数、经济成本、零件的加工工艺特性以及实际使用环境的适应性。经过多番对比，最终确定选用Q235A普通碳素结构钢作为支撑梁材料。该钢材在机械制造和工程结构中应用广泛，焊接性能、塑性和韧性均表现良好，且价格较为低廉，能够满足本设计的承载需求和经济性要求。

其主要性能参数如下表所示：
        \begin{table}[htbp]
        \wuhao
        \centering
        \caption{Q235A号钢的主要性能参数}
        \begin{tabular}{p{3cm}p{4cm}}
        \toprule
        性能指标 & 数值范围/说明 \\
        \midrule
        屈服强度 & $\geq 235 \, \text{MPa}$ \\
        抗拉强度 & $370 \sim 500 \, \text{MPa}$ \\
        延伸率 & $\geq 26\%$ \\
        弹性模量 & $\approx 2.06 \times 10^{5} \, \text{MPa}$ \\
        密度 & $7.85 \, \text{g/cm}^3$ \\
        焊接性能 & 良好 \\
        冲击韧性 & 常温下具有良好韧性 \\
        \bottomrule
        \end{tabular}
        \end{table}
        }
        \subsection{有限元分析}
        \xfbody{材料选定后，对支撑梁开展有限元分析，以检验其在最大受力工况下的强度和安全性。

一是先选择固定位点，如图\ref{fig:fangzhen1}所示。

\begin{figure}[H]
        \centering
        \includegraphics[width=0.5\textwidth]{fangzhen1.png}
        \caption{固定支撑梁的有限元分析示意图}
        \label{fig:fangzhen1}
        \end{figure}
        固定好后，在轮系固定孔处施加1125N的力F。如图\ref{fig:fangzhen2}所示

\begin{figure}[H]
        \centering
        \includegraphics[width=0.5\textwidth]{fangzhen2.png}
        \caption{施加受力的有限元分析示意图}
        \label{fig:fangzhen2}
        \end{figure}
        有限元分析的第一步是对结构模型进行网格划分。网格划分的意图是将连续结构离散化为有限数量的单元，进而借助数值方法求解各节点的应力、应变和位移分布。各单元之间通过节点相互连接，共同反映整体结构的力学特性。如图\ref{fig:fangzhen3}所示，支撑梁被划分为多个有限单元，每个单元的尺寸和形状会根据结构复杂程度及受力情况进行优化，以兼顾分析结果的准确性和计算效率。

\begin{figure}[H]
        \centering
        \includegraphics[width=0.5\textwidth]{fangzhen3.png}
        \caption{支撑梁的网格划分示意图}
        \label{fig:fangzhen3}
        \end{figure}
        网格划分完成后，对模型进行算例运行和受力仿真，利用有限元分析软件计算支撑梁在实际载荷工况下的受力状况，获得结构在不同位置的应力分布。仿真过程中，载荷条件按照AGV自重与额定载荷叠加后施加在支撑梁的受力点上，同时合理设置约束条件和边界条件，以保证计算结果的准确性。如图\ref{fig:fangzhen4}所示，仿真结果呈现了支撑梁在最大受力情况下的应力分布，证实了设计的合理性和安全性。

\begin{figure}[H]
        \centering
        \includegraphics[width=0.5\textwidth]{fangzhen4.png}
        \caption{支撑梁受力仿真结果示意图}
        \label{fig:fangzhen4}
        \end{figure}
        从仿真结果来看，支撑梁的最大应力出现在受力集中区域，将该数值与材料屈服强度进行比对即可判断结构是否满足安全要求。本次仿真结果表明，支撑梁的最大应力为：

\[\sigma_{\mathrm{max}} = 4.891 \times 10^{6} \, \mathrm{N/m^{2}}\]
        远小于 Q235A 钢的屈服强度 235 MPa，说明支撑梁在设计载荷下具有足够的强度和安全余量，能够满足AGV小车在仓储环境中的运行需求。

受力分析和仿真计算完成后，又对结构的位移变形情况进行了模拟评估。从有限元分析结果可以看出，在最大载荷工况下，支撑梁的整体变形量十分微小，模型的最大位移变形为：

\[\delta_{\mathrm{max}} = 1.538 \times 10^{-3} \, \text{mm}\]
        这一数值说明结构在受力作用下的刚度表现良好，变形量远低于一般工程设计中允许的位移限值。换言之，支撑梁在承受设计载荷时几乎没有明显形变，足以保持整体结构的稳定性和精度要求。

从工程视角来看，位移变形的大小直接关系到设备的运行可靠性和使用寿命。倘若变形过大，可能引发零部件之间的干涉或疲劳损伤；而本次仿真结果中的最大位移仅为0.001538 mm，表明设计在刚度和强度方面均满足要求，具备较高的安全余量和使用可靠性。

在受力仿真分析中，除了关注应力与位移的分布情况，还需要对结构的应变进行评估。如图\ref{fig:fangzhen5}和\ref{fig:fangzhen6}所示。

\begin{figure}[H]
        \centering
        \subfigure[支撑梁应变分布示意图]{
        \includegraphics[width=0.5\textwidth]{fangzhen5.png}
        \label{fig:fangzhen5}
        }
        \hfill
        \subfigure[支撑梁最大应变分布示意图]{
        \includegraphics[width=0.5\textwidth]{fangzhen6.png}
        \label{fig:fangzhen6}
        }
        \caption{支撑梁应变分析结果对比图}
        \end{figure}
        应变是衡量材料在外力作用下产生形变程度的重要参数，直接体现了材料的弹性响应能力。由有限元计算结果可知，在最大载荷工况下，结构的最大应变为：

\[\varepsilon_{\mathrm{max}} = 2.268 \times 10^{-5}\]
        这一数值表明支撑梁在受力作用下的变形处于极小范围，远低于材料的屈服应变。换言之，结构在该工况下仍处于弹性阶段，未发生塑性变形，说明设计的刚度和强度均满足要求。

从工程视角来看，最大应变的大小不仅关乎结构安全性，也影响其长期使用性能。较小的应变意味着结构在反复载荷作用下不易产生疲劳损伤，能够维持相对稳定的工作状态。因此本次仿真结果也从侧面表明所选材料与结构设计在理论层面具备合理性。

综合以上分析，通过对AGV小车支撑梁的有限元分析，检验了设计在最大载荷工况下的强度、刚度和安全性：最大应力远低于材料屈服强度，最大位移变形极为微小，最大应变处于弹性范围内，表明设计方案具有充足的安全余量。

}
\chapter{智能仓储移动机器人电气系统设计}
    \section{电气系统总体架构} 
    \xfbody{智能仓储移动机器人的电气系统涵盖电源系统、驱动系统、传感器系统和控制系统四个板块。电源系统为整车各模块提供稳定可靠的电能供应，选用锂离子电池组作为主电源，经由DC-DC变换器为不同电压等级的子系统供电，以满足长时间连续运行的需要；驱动系统由伺服电机和配套驱动电路构成，采用交流伺服电机配合行星减速器，以确保速度和位置控制精度；传感器系统包含红外传感器、超声波传感器和编码器三类，红外传感器负责近距离循迹和避障，超声波传感器负责中远距离测距，电机自带的增量式编码器用于轮速测量和航位推算，为运动控制提供本体状态反馈；控制系统以STM32F103C8T6微控制器为核心，承担传感器数据采集、控制算法执行、驱动信号输出等职责，协调各子系统的工作，实现自主导航和避障功能。以上四个子系统各司其职又紧密协作，共同保障智能仓储移动机器人能够安全高效地完成仓储环境中的物料搬运任务。

\subsection{智能仓储移动机器人主控制器选择}
        \xfbody{在智能仓储移动机器人的电气系统设计中，主控制器的选型是整套方案能否落地的关键环节。综合权衡性能、成本和开发难度等因素后，选定STM32F103C8T6芯片作为主控单元。该芯片基于ARM Cortex-M3内核，运算能力较好、功耗较低、外设接口丰富，能够满足智能仓储移动机器人在导航和避障方面的需求，图\ref{fig:stm32f103c8t6实物图}为其芯片实物图。

\begin{figure}[H]
        \centering
        \includegraphics[width=0.5\textwidth]{stm32f103c8t6实物图.png}
        \caption{STM32F103C8T6芯片实物图}
        \label{fig:stm32f103c8t6实物图}
        \end{figure}
        该控制器主频为72 MHz，配备64 KB程序存储和20 KB的RAM，能够支撑较复杂的算法运行，同时兼容USART、SPI、I2C等多种通信接口，便于与传感器和电机驱动器进行数据交换。加之37个I/O口和12位ADC等资源，开发和扩展都较为便利。

\begin{table}[htbp]
        \wuhao
        \centering
        \caption{STM32F103C8T6芯片参数表}
        \label{tab:stm32f103_compare}
        \begin{tabular}{p{3cm}p{4cm}}
        \toprule
        参数 & 数值 \\
        \midrule
        CPU内核 & ARM Cortex-M3 \\
        CPU最大主频 & 72 MHz \\
        CPU位数 & 32 Bit \\
        程序存储容量 & 64 KB \\
        程序存储器类型 & FLASH \\
        RAM容量 & 20 KB \\
        I/O数量 & 37 \\
        ADC(位数) & 12 bit \\
        \bottomrule
        \end{tabular}
        \end{table}

        图\ref{fig:stm32f103c8t6}为stm32f103c8t6芯片的原理图。

\begin{figure}[H]
        \centering
        \includegraphics[width=0.5\textwidth]{stm32f103c8t6.png}
        \caption{STM32F103C8T6芯片示意图}
        \label{fig:stm32f103c8t6}
        \end{figure}

        }

        }
    \section{电机，传感器等部件选型}
    \xfbody{
        \subsection{电机选型}
        \xfbody{在AGV小车的电机选型过程中，由于本设计采用差速驱动方式，且搬运小车对定位精度和低速平稳性有较高要求，因此驱动电机不仅要提供足够的牵引力，还应具备良好的调速性能和一定的功率裕量。根据设计要求，搬运车最大牵引线速度为30m/min，驱动轮直径为150mm，则驱动轮转速可由式\eqref{eq:4-2-1}计算得到：

        \begin{equation}
        N = \frac{V}{\pi D}
        \label{eq:4-2-1}
        \end{equation}

        其中，$V=30\,\mathrm{m/min}=30000\,\mathrm{mm/min}$，$D=150\,\mathrm{mm}$。为使计算过程中单位保持一致，同时也便于后续与减速箱输出转速对应，此处先将搬运车的最大牵引线速度换算为毫米每分钟，再代入式\eqref{eq:4-2-1}进行计算；驱动轮直径取150mm，既满足车体紧凑布置的要求，又能在一定程度上兼顾通过性与承载能力。代入可得：

        \begin{equation}
        N = \frac{30\times1000}{\pi\times150} \approx 63.69\,\mathrm{r/min}
        \label{eq:4-2-2}
        \end{equation}

        由式\eqref{eq:4-2-2}可知，减速箱输出端对应的转速约为63.69 r/min。进一步分析驱动负载，已知整车自重约200 kg，额定搬运质量为25 kg，满载总质量约225 kg。驱动轮与路面之间的滚动摩擦系数取0.2，按$g\approx10\,\mathrm{m/s^2}$估算，车辆匀速运行时需克服的滚动摩擦力如式\eqref{eq:4-2-3}所示：

        \begin{equation}
        f = 0.2\times225\times10 = 450\,\mathrm{N}
        \label{eq:4-2-3}
        \end{equation}

        由于整车采用双轮差速驱动结构，匀速行驶时两侧驱动轮能够较为均匀地分担整车载荷，故每个驱动轮承担的牵引力可近似视为总牵引力的一半，以便后续计算驱动轮输出转矩和电机功率。单个驱动轮所需牵引力近似为式\eqref{eq:4-2-4}：

        \begin{equation}
        F = \frac{f}{2} = 225\,\mathrm{N}
        \label{eq:4-2-4}
        \end{equation}

        以驱动轮半径$R=75\,\mathrm{mm}=0.075\,\mathrm{m}$进行计算，将单个驱动轮所需承担的牵引力转换为对应的输出转矩。该参数表示驱动电机在减速箱输出端需要提供的扭矩，是保障AGV在满载工况下平稳起步、匀速行驶以及克服滚动阻力的关键依据。单个驱动轮输出转矩如式\eqref{eq:4-2-5}所示：

        \begin{equation}
        M = F\times R = 225\times0.075 = 16.875\,\mathrm{N\cdot m}
        \label{eq:4-2-5}
        \end{equation}

        获得驱动轮输出转矩后，还需进一步估算驱动轮侧的理论机械功率。依据机械功率与转矩、转速之间的对应关系$M=9550\times P_w/n$，将前述转矩值和转速值代入计算，即可得到驱动轮在匀速运行工况下所需的理论机械功率。该结果能够直接反映驱动系统在稳定运行时的能量需求，同时也为后续电机功率校核和减速器匹配提供依据。驱动轮侧的理论机械功率如式\eqref{eq:4-2-6}所示：

        \begin{equation}
        P_w = \frac{M\times n}{9550} = \frac{16.875\times63.69}{9550} \approx 0.1125\,\mathrm{kW}
        \label{eq:4-2-6}
        \end{equation}

        由于实际传动过程中减速箱、联轴器和轴承等部件均会带来一定的机械损耗，理论机械功率无法直接作为电动机的选型依据。若考虑减速箱与传动机构效率$\eta=0.95$，则电动机所需额定功率还需在理论功率基础上加以修正，以确保电机在启动、加速和负载波动等工况下仍具备充足的功率余量。电动机所需额定功率如式\eqref{eq:4-2-7}所示：

        \begin{equation}
        P_m = \frac{P_w}{\eta} = \frac{0.1125}{0.95} \approx 0.1184\,\mathrm{kW}
        \label{eq:4-2-7}
        \end{equation}

        由计算结果可知，理论所需功率约为0.12kW。综合额定功率、输出转矩、控制精度和减速比匹配等因素，最终选用额定功率为400 W的伺服电机SMC80S-0040-30AoK-3DKH作为驱动电机。该电机额定转速为3000 r/min，额定转矩为1.27 N\cdot m，配合减速箱后能够满足驱动轮所需的低速大扭矩输出要求，并为启动、加速和负载波动预留了充足的动力余量；其理论总减速比如式\eqref{eq:4-2-8}所示：

        \begin{equation}
        i = \frac{3000}{63.69} \approx 47.1
        \label{eq:4-2-8}
        \end{equation}

        需要说明的是，上述选型计算主要基于匀速行驶时的滚动阻力工况。实际运行中，AGV还面临以下动态载荷情形：①起步与制动阶段的惯性加速力，以$0.3\,\mathrm{m/s^2}$加速度估算，附加惯性力约$225\times0.3=67.5\,\mathrm{N}$，折算到单轮约增加转矩$2.5\,\mathrm{N\cdot m}$；②差速转弯时外侧轮承载增大，极端情况下载荷偏移约10\%$\sim$15\%；③举升过程中搬运物抬升导致重心上移，对驱动轮法向力的影响有限但需考虑；④地面不平整引起的附加载荷波动。综合上述工况，实际峰值转矩需求约为稳态值的1.3$\sim$1.5倍，即约$22\sim25\,\mathrm{N\cdot m}$。所选电机配合47.1倍减速箱后的峰值输出转矩为$1.27\times3\times47.1\approx179.9\,\mathrm{N\cdot m}$（伺服电机过载系数一般为3倍额定转矩），远高于动态峰值需求，具备充足的动力余量。因此，所选伺服电机能够满足搬运车30 m/min的最高行驶速度要求以及各类工况下的牵引需求。相关型号参数如表格\ref{tab:smc_servo}所示。

\begin{table}[htbp]
        \wuhao
        \centering
        \caption{SMC系列伺服电机技术参数}
        \label{tab:smc_servo}
        \begin{tabular}{C{1.9cm}C{2.7cm}C{2.5cm}C{2.5cm}C{2.5cm}C{2.5cm}}
        \toprule
        \multicolumn{2}{c}{伺服电机型号} & \tabincell{c}{SMC60S-\\0020-30\\AoK-3DKH} & \tabincell{c}{SMC60S-\\0040-30\\AoK-3DKH} & \tabincell{c}{SMC80S-\\0040-30\\AoK-3DKH} & \tabincell{c}{SMC80S-0075-30\\AoK-3DKH} \\
        \midrule
        \multicolumn{2}{c}{适配驱动} & \tabincell{c}{FD123-LA-000\\FD123-CA-000\\FD123-CC-000\\FD123-EA-000} & \tabincell{c}{FD123-LA-000\\FD123-CA-000\\FD123-CC-000\\FD123-EA-000} & \tabincell{c}{FD133-LA-000\\FD133-CA-000\\FD133-CC-000\\FD133-EA-000} & \tabincell{c}{FD133-LA-000\\FD133-CA-000\\FD133-CC-000\\FD133-EA-000} \\
        \multicolumn{2}{c}{直流母线电压 $U_{DC}$} & 48 & 48 & 48 & 48 \\
        \midrule
        \multirow{4}{*}{\tabincell{c}{连续\\特性}} & \tabincell{c}{额定功率 $P_n$(W)} & 200 & 400 & 400 & 750 \\
        & \tabincell{c}{额定转矩 $T_n$(Nm)} & 0.64 & 1.27 & 1.27 & 2.39 \\
        & \tabincell{c}{额定转速 $n_N$(rpm)} & 3000 & 3000 & 3000 & 3000 \\
        & \tabincell{c}{额定电流 $I_n$(A)} & 4.6 & 10 & 9.6 & 21.5 \\
        \midrule
        & \tabincell{c}{瞬时最大转矩 $T_m$(Nm)} & 1.92 & 3.81 & 3.81 & 7.17 \\
        & \tabincell{c}{瞬时最大电流 $I_m$(A)} & 13.8 & 25 & 24 & 64.5 \\
        & \tabincell{c}{连续静态转矩 $T_s$(Nm)} & 0.7 & 1.4 & 1.4 & 2.63 \\
        & \tabincell{c}{连续静态电流 $I_s$(A)} & 5.06 & 11 & 10.6 & 23.7 \\
        & \tabincell{c}{线-线电阻 $R_L$\\($\Omega$)} & 1.1 & 0.42 & 0.22 & 0.1 \\
        & \tabincell{c}{线-线电感 $L_L$(mH)} & 2.4 & 0.79 & 1 & 0.46 \\
        & \tabincell{c}{机械时间常数(ms)} & 3.22 & 1.84 & 1.65 & 1.79 \\
        & \tabincell{c}{反电势常数 $K_e$(V/krpm)} & 9 & 8 & 8 & 6.7 \\
        & \multicolumn{1}{c}{位置反馈装置} & \multicolumn{4}{c}{2500ppr 光电式编码器（增量式）} \\
        & \multicolumn{1}{c}{冷却方式} & \multicolumn{4}{c}{全封闭，自冷却} \\
        \bottomrule
        \end{tabular}
        \end{table}

        所选SMC80S-0040-30AoK-3DKH伺服电机满足本设计的功率与转矩需求。其电机外形图如图\ref{fig:smc80s}所示。

\begin{figure}[H]
        \centering
        \includegraphics[width=0.5\textwidth]{SMC80S-0040-30AoK-3DKH电机外形图.png}
        \caption{SMC80S-0040-30AoK-3DKH伺服电机外形图}
        \label{fig:smc80s}
        \end{figure}
        }

        \subsection{减速器选型}
        \xfbody{
        依据前述计算结果，驱动电机的额定转速为$n_1=3000\,\mathrm{r/min}$，驱动轮在满足整车最高行驶速度要求时对应的转速为$n_2=63.69\,\mathrm{r/min}$。据此，电机输出端与驱动轮之间所需的理论总传动比可按式：

        \begin{equation}
        i=\frac{n_1}{n_2}=\frac{3000}{63.69}\approx47.10
        \end{equation}

        进行计算。不难看出，要将电机输出转速降至驱动轮所需范围，同时保证输出扭矩满足满载运行要求，必须配置具有较大减速比的减速器。考虑到AGV小车本身结构紧凑、安装空间有限，且对传动平稳性、输出精度和维护便利性要求较高，故不宜采用多级串联式减速方案，而应优先选用体积较小、传动效率较高、减速比范围较宽的标准化减速器。

        综合传动比匹配、安装布置和产品系列的可获得性，最终选定ZJPX系列精密行星减速机作为电机与驱动轮之间的减速装置。相较于普通减速方案，这类减速机结构更紧凑、传动更平稳，也更适合安装空间有限的AGV平台。其主要优点如下：
        \begin{enumerate}
        \item 采用齿轮啮合传动，输出扭矩较大，传动效率高，且运行噪声小，适合对平稳性要求较高的AGV系统；
        \item 传动过程中受力分布较均匀，能够改善整机运行的稳定性，减少冲击和振动对机构的影响；
        \item 结构紧凑，输入端与输出端的连接方式灵活，便于与伺服电机组合安装，整体装配可靠性高；
        \item 日常维护简单，润滑条件较稳定，能够满足AGV长期连续运行的使用需求。
        \end{enumerate}

        ZJPX系列精密行星减速机的主要技术参数如表\ref{tab:zjpx_param}所示。
        \begin{table}[htbp]
        \wuhao
        \centering
        \caption{ZJPX行星减速机技术参数}
        \label{tab:zjpx_param}
        \begin{tabular}{C{1.6cm}C{1.75cm}C{1.75cm}C{1.75cm}C{1.75cm}C{1.2cm}C{1.2cm}}
        \toprule
        \tabincell{c}{产品\\型号} & ZJPX65 & ZJPX85 & ZJPX115 & ZJPX142 & 单位 & 级数 \\
        \midrule
        \tabincell{c}{最大\\径向力} & 1550 & 3055 & 4330 & 9480 & N & \multirow{2}{*}{} \\
        \tabincell{c}{最大\\轴向力} & 1220 & 2330 & 3300 & 6800 & N & \\
        \midrule
        \multirow{3}{*}{\tabincell{c}{满载\\效率}} & \multicolumn{4}{c}{95} & \multirow{3}{*}{\%} & 1级 \\
        & \multicolumn{4}{c}{93} & & 2级 \\
        & \multicolumn{4}{c}{90} & & 3级 \\
        \midrule
        \tabincell{c}{平均\\寿命} & \multicolumn{4}{c}{20000} & h & \multicolumn{1}{c}{} \\
        \bottomrule
        \end{tabular}
        \end{table}

        根据减速器资料可知，驱动电机经过减速箱后输出扭矩可达到90 N·m。结合整车的安装空间、传动效率以及运行稳定性要求，最终确定采用ZJPX115型二级精密行星减速器，其减速比为40。该减速比与理论值47.10存在差异，工程实现中通过伺服驱动器限速进行补偿：当电机最高转速限制在约$2550\,\mathrm{r/min}$时，减速后轮速约$63.75\,\mathrm{r/min}$，对应线速度约$30\,\mathrm{m/min}$，可满足速度指标。该方案既能够满足驱动轮所需的输出转矩，又能够保证整车在低速运行时具有较好的控制精度和运动平稳性。

        }
        \subsection{传感器选型}
        \xfbody{
        AGV配备红外传感器和超声波传感器两种类型，二者在检测距离、响应速度和适用场景方面各具优势，通过配合使用可实现互补。红外传感器基于光电反射原理，适合近距离检测，具有响应速度快、体积小、成本低的特点，常用于地面循迹线检测；超声波传感器基于声波飞行时间测距原理，适合中远距离障碍物检测，精度较高且不受光照条件影响。两类传感器协同使用，分别承担循迹和避障功能，使AGV对周围环境的感知更加全面。

红外线传感器选用型号为TCRT5000的反射式红外传感器，如图\ref{fig:TCRT5000}所示。

\begin{figure}[H]
        \centering
        \includegraphics[width=0.5\textwidth]{TCRT5000.jpg}
        \caption{TCRT5000红外传感器实物图}
        \label{fig:TCRT5000}
        \end{figure}TCRT5000是一种集发射与接收于一体的反射式光电传感器，由红外发光二极管和NPN光电三极管组成。其工作原理为：红外发射管向前方发射红外光，当遇到障碍物或地面黑色引导线时，红外光被反射回来并由接收管检测，依据反射光强度判断前方是否存在检测目标。本文选用带比较器电路的TCRT5000传感器模块，该模块板载LM393比较器和电位器，可输出数字量（高低电平），直接连接STM32的GPIO引脚即可使用。该模块的有效检测距离为$1\sim25\ \text{mm}$，适合黑白线循迹检测；但对AGV避障而言，$25\ \text{mm}$的检测距离不足以提供安全制动裕量，因此本文中TCRT5000主要用于地面轨迹线检测（循迹），障碍物检测功能由测距范围更远的HC-SR04超声波传感器承担。

        超声波传感器选用型号为HC-SR04的超声波测距模块，如图\ref{fig:HC-SR04}所示。

\begin{figure}[H]
        \centering
        \includegraphics[width=0.5\textwidth]{HC-SR04.jpg}
        \caption{HC-SR04超声波传感器实物图}
        \label{fig:HC-SR04}
        \end{figure}
        该模块通过发射超声波脉冲并接收障碍物反射信号来测量距离。其工作过程为：STM32向Trig引脚发送一个不低于$10\,\mu\mathrm{s}$的高电平触发信号，模块随即发射一组40kHz的超声波脉冲，同时将Echo引脚置为高电平；当超声波遇到障碍物反射回来被接收后，Echo引脚恢复低电平，通过测量Echo高电平的持续时间即可计算出障碍物距离。HC-SR04工作电压为5V，测距范围为2cm至400cm，测量精度可达3mm，适合AGV在仓储环境中进行中远距离障碍物检测。将该模块与STM32连接后，AGV能够实时获取前方障碍物的距离信息，与红外传感器的近距离检测相配合，使导航和避障功能更加完善。

        两种传感器协同工作，覆盖了AGV在不同距离范围的感知需求，为后续运动控制提供了可靠的数据支撑。

除环境感知传感器外，AGV还需要获取自身运动状态。所选SMC80S伺服电机自带增量式光电编码器，分辨率为$2500\,\text{P/R}$，经4倍频后每转输出$10000$个脉冲，可精确测量左右轮的转速和转角。通过对编码器脉冲进行计数和微分运算，能够实时计算出各轮的角速度和累计转角，进而由差速运动学方程推算出机器人的航向角变化量$\Delta\theta$和行驶里程，构成航位推算（Dead Reckoning）系统。编码器信号通过STM32的定时器编码器接口（TIMx\_CH1/TIMx\_CH2）以四倍频模式采集，无需额外的信号调理电路。该方案在短距离、平坦地面条件下能够满足定位精度需求；若后续需要更高精度，可引入惯性测量单元（IMU）进行航向角融合修正。

}
        }
    \section{控制电路设计}
        \subsection{STM32F103C8T6最小系统}
        \xfbody{STM32F103C8T6最小系统主要由电源电路、滤波电路、时钟电路、SWD程序下载调试电路和复位电路组成。如图\ref{fig:stm32f103c8t6_min}所示。

\begin{figure}[H]
        \centering
        \includegraphics[width=0.8\textwidth]{最小系统原理图.png}
        \caption{STM32F103C8T6最小系统示意图}
        \label{fig:stm32f103c8t6_min}
        \end{figure}
        其中，电源电路采用ME6211C33M5G-N。选用ME6211C33芯片，如图\ref{fig:ME6211C33}所示。

        \begin{figure}[H]
        \centering
        \includegraphics[width=0.5\textwidth]{电源.png}
        \caption{ME6211C33芯片实物图}
        \label{fig:ME6211C33}
        \end{figure}

        该芯片采用SOT-23-5封装，是一款低功耗低压差LDO芯片，工作电压范围2V~6V，最大输出电流500mA，待机电流仅0.1$\mu$A，具备过热保护和过流保护功能，可为STM32提供稳定的3.3V供电。

滤波电路主要由电容器构成，其功能是稳定电源电压。如图\ref{fig:power_filter}所示为电源滤波电路。

        \begin{figure}[H]
        \centering
        \includegraphics[width=0.5\textwidth]{电源滤波.png}
        \caption{电源滤波电路示意图}
        \label{fig:power_filter}
        \end{figure}

        为保障STM32供电电压的稳定性，在电源正负极之间需要放置一个104去耦电容，即$10\times10^4\,\mathrm{pF}$（100 nF），利用电容隔直通交的特性进行去耦，滤除电路中的电压波动和高频振荡，降低电源噪声，使电源输出更加平稳。

时钟电路采用外部晶振，频率为8MHz。如图\ref{fig:crystal_oscillator}所示为晶振电路示意图。

\begin{figure}[H]
        \centering
        \includegraphics[width=0.5\textwidth]{时钟电路.png}
        \caption{时钟电路示意图}
        \label{fig:crystal_oscillator}
        \end{figure}
        选用8MHz晶振可为STM32F103C8T6控制器提供稳定的时钟信号。外部晶振相较于内部RC振荡器具有更高的频率稳定性和精度。为满足晶振的起振要求，晶振电路还需连接两个12pF的电容。

SWD程序下载调试电路用于实现STM32F103C8T6的程序烧录和调试功能，确保开发过程中能够便捷地进行代码测试和问题排查。如图\ref{fig:swd_circuit}所示为SWD程序下载调试电路示意图。

\begin{figure}[H]
        \centering
        \includegraphics[width=0.5\textwidth]{swd.png}
        \caption{SWD程序下载调试电路示意图}
        \label{fig:swd_circuit}
        \end{figure}
        SWD（Serial Wire Debug）是一种两线调试协议，由SWDIO（数据线）和SWCLK（时钟线）组成。借助SWD接口，开发者可以使用ST-Link等调试器与STM32通信，完成程序烧录、断点调试和寄存器监视等操作。

复位电路的作用是确保STM32F103C8T6控制器在系统异常或电源波动等情况下能够及时复位，恢复正常运行状态。如图\ref{fig:reset_circuit}所示为复位电路示意图。

\begin{figure}[H]
        \centering
        \includegraphics[width=0.5\textwidth]{reset.png}
        \caption{复位电路示意图}
        \label{fig:reset_circuit}
        \end{figure}
        复位电路在系统异常时将MCU恢复到初始状态。STM32的复位引脚（NRST）内部设有上拉电阻，默认为高电平，拉低时触发复位。电路由按键、电容和上拉电阻组成：按键按下时NRST被拉低，触发复位；电容用于滤除按键抖动；上拉电阻保证未按下时引脚为高电平，防止误触发。

        STM32F103C8T6最小系统的设计涵盖了电源、滤波、时钟、SWD调试和复位等关键电路，为AGV小车的运动控制和环境感知奠定了硬件基础。

        }
        \subsection{传感器接口电路设计}
        \xfbody{
        传感器接口电路涵盖红外传感器接口和超声波传感器接口两部分，保障AGV能够准确获取环境信息。

        红外传感器接口电路用于连接TCRT5000，保障其与STM32的正常通信。如图\ref{fig:ir_interface}所示。

        \begin{figure}[H]
        \centering
        \includegraphics[width=0.5\textwidth]{ir_interface.png}
        \caption{红外线传感器接口电路示意图}
        \label{fig:ir_interface}
        \end{figure}
        TCRT5000输出为数字量，直接连接STM32的GPIO引脚即可。信号线上加10k$\Omega$上拉电阻拉高到3.3V，防止信号浮空；再加100nF去耦电容滤除高频噪声，增强检测的稳定性。

        超声波传感器接口电路用于连接HC-SR04模块。如图\ref{fig:ultrasonic_interface}所示。

        \begin{figure}[H]
        \centering
        \includegraphics[width=0.5\textwidth]{ultrasonic_interface.png}
        \caption{超声波传感器接口电路示意图}
        \label{fig:ultrasonic_interface}
        \end{figure}
        HC-SR04模块设有四个引脚：VCC、GND、Trig和Echo。VCC和GND分别连接到5V电源和地线，确保模块正常供电；Trig引脚连接到STM32F103C8T6的PA1输出引脚，用于触发超声波发射；Echo引脚连接到STM32F103C8T6的PA2输入引脚，用于接收反射信号。由于Echo为5V电平，实际硬件实现需确认所用引脚的5V容限条件，必要时应增加分压或电平转换电路，以保证接口可靠性。

        红外传感器接口电路负责地面循迹线检测，超声波传感器接口电路负责中远距离障碍物测距，二者相互配合，为循迹控制和避障决策提供数据支撑。

}


\chapter{智能仓储移动机器人运动学建模与仿真}
    \section{差速驱动底盘运动学方程推导}
    \xfbody{在对智能仓储移动机器人的运动性能进行分析和控制之前，需要建立能够反映其结构特性和运动约束的运动学模型。运动学模型描述的是机器人位姿（位置与姿态）随时间的变化规律，不涉及力和力矩等动力学因素，是开展轨迹规划、速度控制和路径跟踪的基础。对于本文所研究的双轮差速智能仓储移动机器人，其运动主要发生在二维平面内，故可借助平面直角坐标和航向角建立运动学方程，为后续的轨迹仿真和控制器设计提供数学基础。

    所研究的智能仓储移动机器人为双轮差速移动机器人，模型如图\ref{双轮模型.png}所示，下面展开运动学分析。

    \begin{figure}[H]
    \centering
    \includegraphics[width=0.5\textwidth]{双轮模型.png}
    \caption{双轮差速移动机器人模型示意图}
    \label{双轮模型.png}
    \end{figure}

    设机器人在二维平面内运动，定义全局坐标系$O-XY$和机器人坐标系$O'-X'Y'$，其中$O'$为机器人质心位置。机器人具有两个驱动轮，分别为左轮和右轮，轮距为$L$，轮半径为$r$。设左轮和右轮的线速度分别为$v_L$和$v_R$，则机器人的线速度$v$和角速度$\omega$可以表示为：

    \begin{equation}
    v = \frac{v_R + v_L}{2}
    \label{eq:linear_velocity}
    \end{equation}
    \begin{equation}
    \omega = \frac{v_R - v_L}{L}
    \label{eq:angular_velocity}
    \end{equation}

    机器人在全局坐标系中的位置$(x, y)$和姿态$\theta$的变化率可由以下运动学方程描述：
    
    \begin{equation}
    \dot{x} = v \cos\theta
    \label{eq:dot_x}
    \end{equation}
    \begin{equation}
    \dot{y} = v \sin\theta
    \label{eq:dot_y}
    \end{equation}
    \begin{equation}\dot{\theta} = \omega
    \label{eq:dot_theta}
    \end{equation}

    为便于在MATLAB中直接以线速度和角速度形式构造典型轨迹，进一步引入轮角速度作为执行层输入。由差速底盘几何关系可得右轮和左轮的角速度分别为：

    \begin{equation}
    \omega_R = \frac{2v + \omega L}{2r}
    \label{eq:right_wheel_speed}
    \end{equation}
    \begin{equation}
    \omega_L = \frac{2v - \omega L}{2r}
    \label{eq:left_wheel_speed}
    \end{equation}

    上述关系式与式\eqref{eq:dot_x}至式\eqref{eq:dot_theta}共同构成差速机器人典型轨迹仿真的基础，只需给定不同的$v(t)$和$\omega(t)$，即可通过积分得到机器人在平面内的位姿变化，并分析轨迹长度、位移、曲率及左右轮角速度变化规律。
    }
    \section{典型运动轨迹仿真与分析}
    \xfbody{依据附录程序\texttt{chasu.m}，在MATLAB环境下构建了差速机器人典型轨迹运动学正解仿真模型。该仿真不涉及控制器，而是直接给定线速度指令$v(t)$和角速度指令$\omega(t)$，通过前向欧拉积分方法对式\eqref{eq:dot_x}至式\eqref{eq:dot_theta}进行离散求解，获得机器人在平面内的位姿变化。仿真参数为：机器人初始位姿设为$(0,0,0)$，轮半径取$r=0.075\ \text{m}$，轮距取$L=0.52\ \text{m}$，积分步长为$0.01\ \text{s}$，总仿真时间为$20\ \text{s}$。在此基础上，通过构造不同的速度指令组合，分别得到直线轨迹、定曲率圆轨迹、S形轨迹和8字轨迹四类典型运动形式，用以分析差速底盘在不同工况下的轨迹特征和轮速变化规律。

1. 直线运动是最基本的运动形式，其控制输入设定为恒定线速度$v=0.45\ \text{m/s}$、角速度$\omega=0$。其运动轨迹如图\ref{fig:straight_line}所示。

\begin{figure}[H]
    \centering
    \includegraphics[width=0.5\textwidth]{直线.png}
    \caption{直线运动轨迹仿真结果}
    \label{fig:straight_line}
    \end{figure}
    由仿真结果可知，机器人沿$x$轴正方向做匀速直线运动，20 s内累计路径长度约为$9.0\ \text{m}$，末端位移同样为$9.0\ \text{m}$，终点航向角保持为$0^\circ$。左右轮角速度始终相等，表明此时差速底盘处于标准直行状态。

2. 圆弧运动对应仓储AGV常见的连续转向工况。程序中设置$v=0.35\ \text{m/s}$、$\omega=0.45\ \text{rad/s}$，即线速度和角速度均保持常值，由此形成定曲率圆轨迹。其运动轨迹如图\ref{fig:circular_arc}所示。

仿真结果表明，机器人在20 s内完成长度约为$7.0\ \text{m}$的圆弧路径，终点位移约为$1.521\ \text{m}$，终点航向角约为$155.662^\circ$。由于$\omega/v$保持常数，轨迹曲率稳定不变，左右轮角速度始终存在固定差值，体现了差速驱动借助恒定速差实现平滑转弯的特点。

\begin{figure}[H]
    \centering
    \includegraphics[width=0.32\textwidth]{圆形.png}
    \caption{圆弧运动轨迹仿真结果}
    \label{fig:circular_arc}
    \end{figure}

    3. S形运动用来反映机器人在左右交替转向场景下的运动状态。程序中给定恒定线速度$v=0.40\ \text{m/s}$，并令角速度满足$\omega(t)=0.6\sin(0.5t)$，使机器人在前进过程中周期性改变转向方向。其运动轨迹如图\ref{fig:s_shape}所示。

\begin{figure}[H]
    \centering
    \includegraphics[width=0.5\textwidth]{s型.png}
    \caption{S形运动轨迹仿真结果}
    \label{fig:s_shape}
    \end{figure}
    仿真结果显示，该轨迹总路径长度约为$8.0\ \text{m}$，终点位移约为$5.269\ \text{m}$，最大角速度达到$0.6\ \text{rad/s}$，最大曲率约为$1.5\ \text{m}^{-1}$。从轨迹形态来看，机器人在整体前进的同时完成了多次连续换向，较好地体现了差速底盘在曲折通道中的灵活性。

4. 8字形运动进一步增加了转向频率和曲率变化幅值。程序中设置线速度$v=0.38\ \text{m/s}$，角速度为$\omega(t)=0.7\sin(0.7t)$，由此形成具有明显自交特征的8字形轨迹。其运动轨迹如图\ref{fig:figure_eight}所示。

\begin{figure}[H]
    \centering
    \includegraphics[width=0.5\textwidth]{8字.png}
    \caption{8字形运动轨迹仿真结果}
    \label{fig:figure_eight}
    \end{figure}
    从结果可以看出，8字轨迹总路径长度约为$7.6\ \text{m}$，终点位移约为$5.817\ \text{m}$，最大角速度达到$0.7\ \text{rad/s}$，最大曲率约为$1.842\ \text{m}^{-1}$。与S形轨迹相比，该工况对机器人姿态变化速度和轮速协调能力的要求更高，更适合用于验证差速底盘在复杂曲线路径下的运动学响应特性。

四类轨迹的输入指令与左右轮角速度变化曲线如图\ref{fig:control_input}所示。从图中可以看出，直线轨迹下左右轮角速度完全重合；定曲率圆轨迹下两轮角速度保持恒定差值；S形和8字轨迹下左右轮角速度均随时间呈周期性变化，且8字轨迹的轮速波动幅值更大，表明其对差速执行机构的动态调节能力要求更高。

\begin{figure}[H]
    \centering
    \includegraphics[width=0.5\textwidth]{输入与轮速曲线.png}
    \caption{控制输入与轮速曲线}
    \label{fig:control_input}
    \end{figure}
    各类轨迹的瞬时曲率对比如图\ref{fig:curvature_comparison}所示。依据$\kappa=\omega/v$可知，直线轨迹曲率始终为0，定曲率圆轨迹曲率保持常值，而S形轨迹和8字轨迹的曲率均呈正负交替变化，其中8字轨迹的峰值更大，表明其路径弯曲程度更高。

\begin{figure}[H]
    \centering
    \includegraphics[width=0.5\textwidth]{曲率变化.png}
    \caption{曲率对比图}
    \label{fig:curvature_comparison}
    \end{figure}
    为了更直观地对比不同轨迹的运动学特征，将仿真结果汇总如表\ref{tab:traj_summary}所示。

\begin{table}[H]
    \centering
    \caption{典型轨迹仿真结果汇总}
    \label{tab:traj_summary}
    \begin{tabular}{p{2.5cm}p{1.8cm}p{1.5cm}p{2cm}p{2.2cm}p{2cm}}
    \toprule
    轨迹类型 & 路径长度/m & 位移/m & 平均速度/(m/s) & 最大角速度/(rad/s) & 最大曲率/(m$^{-1}$) \\
    \midrule
    直线轨迹 & 9.0 & 9.0 & 0.45 & 0 & 0 \\
    定曲率圆轨迹 & 7.0 & 1.521 & 0.35 & 0.45 & 1.286 \\
    S形轨迹 & 8.0 & 5.269 & 0.40 & 0.60 & 1.500 \\
    8字轨迹 & 7.6 & 5.817 & 0.38 & 0.70 & 1.842 \\
    \bottomrule
    \end{tabular}
    \end{table}

    以上仿真结果显示，基于\texttt{chasu.m}的差速机器人运动学仿真能够反映不同速度指令组合下的轨迹变化规律：直线轨迹体现了底盘的基础直行能力，定曲率圆轨迹反映了稳态转弯特性，S形和8字轨迹则呈现了机器人在连续变曲率工况下的姿态调整能力。
    }
    \section{路径规划与循迹仿真}
    \xfbody{在完成差速底盘运动学建模和典型轨迹仿真验证的基础上，进一步针对仓储搬运场景开展路径规划与循迹控制仿真。本节采用两层导航架构：全局路径规划层采用Hybrid A*算法在包含货架障碍物的栅格地图上离线规划无碰撞路径，确定AGV从起点到终点的几何轨迹；局部循迹跟踪层以规划路径为参考线铺设引导标识，模拟5路TCRT5000红外传感器检测轨迹偏移量，通过PID控制器进行差速纠偏，实现在线循迹跟踪。需要说明的是，Hybrid A*路径规划解决的是"走哪条路"的问题，而红外循迹解决的是"如何沿路走"的问题，二者分别对应导航系统的决策层和执行层。采用Hybrid A*而非直接铺设固定磁条或二维码路线的原因在于：仓储货架布局可能随业务需求调整，Hybrid A*能够在地图变更后自动重新规划路径，无需人工重新铺设引导标识；同时，Hybrid A*生成的路径满足差速底盘的转弯半径约束，保证底层控制器可直接执行，而人工铺设的路径未必满足运动学可行性。此外，两层架构具有良好的可扩展性——当后续引入激光雷达或视觉SLAM等更先进的感知手段时，只需替换规划层，循迹执行层可保持不变。

    \subsection{Hybrid A*路径规划算法}
    \xfbody{Hybrid A*算法在传统A*搜索算法基础上加入了非完整约束的运动学可行性扩展，适用于差速驱动等具有转弯半径限制的移动机器人。与标准A*在离散栅格上仅允许8邻域移动不同，Hybrid A*的节点状态包含位置和航向角$(x, y, \theta)$，扩展时采用符合差速运动学的运动基元生成后继节点，从而保证规划路径可直接被底层控制器执行。

    本文采用的运动基元基于差速驱动模型。由式\eqref{eq:linear_velocity}和式\eqref{eq:angular_velocity}可知，差速底盘的线速度$v=(v_R+v_L)/2$、角速度$\omega=(v_R-v_L)/L$，故等效转弯半径为$R = v/\omega = L(v_R+v_L)/[2(v_R-v_L)]$，对应等效曲率$\kappa = 1/R$。在路径规划中，选取5种离散曲率$\kappa \in \{-1.11, -0.52, 0, 0.52, 1.11\}\ \text{m}^{-1}$（分别对应转弯半径$R$为$0.90\ \text{m}$、$1.94\ \text{m}$、$\infty$、$1.94\ \text{m}$、$0.90\ \text{m}$，等效左右轮速度比约为$1.8:1$、$1.3:1$、$1:1$、$1:1.3$、$1:1.8$）和前进/后退两种运动方向，共生成10种运动基元。每个基元的步长取$d=0.5\ \text{m}$，经运动学积分得到新节点位姿：

    \begin{equation}
    \begin{cases}
    x' = x + d \cdot \cos(\theta + \Delta\theta/2) \\
    y' = y + d \cdot \sin(\theta + \Delta\theta/2) \\
    \theta' = \theta + \Delta\theta
    \end{cases}
    \label{eq:motion_primitive}
    \end{equation}

    其中$\Delta\theta = d \cdot \kappa$为航向变化量。该运动模型无需转向机构，航向变化完全由左右轮速差产生，与差速底盘的驱动原理一致。代价函数采用$ f = g + h $的形式，其中$g$为从起点到当前节点的实际代价（含转弯惩罚和倒车惩罚），$h$为当前节点到目标点的欧氏距离启发值。

    仿真地图设置为$10\ \text{m} \times 8\ \text{m}$的仓储环境，栅格分辨率为$0.2\ \text{m}$，包含5个矩形货架障碍物。起点设为$(1.0, 1.0)$，航向$0^\circ$；终点设为$(9.0, 7.0)$，航向$90^\circ$。

    }

    \subsection{路径规划仿真结果}
    \xfbody{基于上述参数运行Hybrid A*算法，搜索过程在第3134次迭代时成功找到可行路径，生成的规划路径如图\ref{fig:hybrid_astar_result}所示。

    \begin{figure}[H]
        \centering
        \includegraphics[width=0.55\textwidth]{hybrid_astar_result.png}
        \caption{Hybrid A*路径规划结果}
        \label{fig:hybrid_astar_result}
    \end{figure}

    从图中可以看出，规划路径成功绕过了所有货架障碍物，从起点到达终点。路径总长度约为$10.71\ \text{m}$，共包含24个路径节点，平均节点间距约为$0.48\ \text{m}$。路径整体较为平滑，在障碍物间隙处能够合理选择转弯方向，没有明显的锯齿或多余绕行。与标准A*算法相比，Hybrid A*的运动基元基于差速运动学生成，规划路径天然满足转弯半径约束，可直接被底层控制器执行；而标准A*在离散栅格上仅允许8邻域移动，规划路径存在锯齿状折线，必须经过后处理才能用于实际机器人。在搜索完成后，对规划路径施加了一次轻量级的均值平滑（3次迭代），以消除节点间距不均匀导致的微小折角，平滑操作不改变路径的运动学可行性。

    为满足后续循迹仿真的精度要求，对原始规划路径进行等间距插值处理，插值间距为$0.05\ \text{m}$，插值后路径点数增加至215个，路径总长度约为$10.71\ \text{m}$。

    }

    \subsection{红外传感器循迹模型}
    \xfbody{获得参考路径后，采用5路TCRT5000反射式红外传感器阵列进行轨迹线检测，实现循迹控制。传感器阵列安装在车体前端、驱动轮轴线前方$80\ \text{mm}$处，5个传感器等间距布置，间距为$40\ \text{mm}$，各传感器相对车体中心线的横向偏移量分别为$-80\ \text{mm}$、$-40\ \text{mm}$、$0\ \text{mm}$、$40\ \text{mm}$、$80\ \text{mm}$，阵列总覆盖宽度为$160\ \text{mm}$，如图\ref{fig:sensor_layout}所示。

    \begin{figure}[H]
        \centering
        \includegraphics[width=0.35\textwidth]{sensor_layout.png}
        \caption{5路红外传感器布局示意图}
        \label{fig:sensor_layout}
    \end{figure}

    当传感器检测到黑线（反射率低于阈值）时输出为1，否则输出为0。依据5路传感器的检测状态，按照权重$\{-4, -2, 0, +2, +4\}$计算加权平均偏移量$e$，反映车体相对于轨迹线的横向偏差：

    \begin{equation}
    e = \frac{\sum_{i=1}^{5} w_i \cdot s_i}{\sum_{i=1}^{5} s_i}
    \label{eq:sensor_offset}
    \end{equation}

    其中$w_i$为第$i$个传感器的权重，$s_i \in \{0, 1\}$为检测状态。偏移量$e$的取值范围为$[-2, +2]$，负值表示偏左，正值表示偏右。当所有传感器均未检测到黑线时，说明车体横向偏移已超出传感器阵列的检测范围，此时将上一次的偏移量乘以1.5倍系数进行外推，使控制器产生更强的纠偏力矩，引导车体尽快回到可检测区域内，避免因持续丢线导致的控制失效。

    循迹控制器采用PID算法，以偏移量$e$为输入，输出角速度指令$\omega$，再经由差速运动学转换为左右轮转速差，实现纠偏：

    \begin{equation}
    \omega = K_p \cdot e + K_i \cdot \int e\,\mathrm{d}t + K_d \cdot \frac{\mathrm{d}e}{\mathrm{d}t}
    \label{eq:line_pid}
    \end{equation}

    本文取$K_p=3.0$、$K_i=0.08$、$K_d=1.0$，角速度限幅为$\pm 1.5\ \text{rad/s}$，AGV标称行驶速度为$0.3\ \text{m/s}$。

    需要说明的是，此处的PID控制器用于将红外传感器阵列的横向偏移量映射为角速度指令，属于单层循迹控制方案，目的在于验证全局路径规划与局部引导线跟踪协同工作的可行性。关于双环级联PID（位姿外环+轮速内环）的系统设计与性能验证，将在第六章详细展开。

    }

    \subsection{循迹仿真结果与分析}
    \xfbody{将Hybrid A*规划并插值后的路径作为参考线，运行PID循迹仿真，结果如图\ref{fig:line_following_result}所示。

    \begin{figure}[H]
        \centering
        \includegraphics[width=0.55\textwidth]{line_following_result.png}
        \caption{PID循迹仿真结果}
        \label{fig:line_following_result}
    \end{figure}

    从轨迹对比图可以看出，循迹轨迹与参考路径基本吻合，在转弯区段存在小幅偏移，但整体跟踪效果良好。横向误差曲线显示，误差在起步阶段较大，随后快速稳定并保持在较小范围内波动。

    循迹仿真结果汇总如表\ref{tab:line_following_result}所示。

    \begin{table}[H]
        \centering
        \caption{PID循迹仿真结果}
        \label{tab:line_following_result}
        \begin{tabular}{ccc}
        \toprule
        RMS横向误差/m & 最大横向误差/m & 仿真时长/s \\
        \midrule
        0.008 & 0.016 & 35.1 \\
        \bottomrule
        \end{tabular}
    \end{table}

    仿真结果表明，所设计的5路红外传感器PID循迹方案能够在仓储环境下较好地跟踪Hybrid A*规划路径，RMS横向误差为$0.8\ \text{cm}$，最大横向误差为$1.6\ \text{cm}$。阵列总覆盖宽度为$160\ \text{mm}$（$\pm 80\ \text{mm}$），最大横向误差远在传感器覆盖范围内，表明传感器布局具有充足的检测裕量。循迹过程耗时$35.1\ \text{s}$，控制器响应及时，纠偏过程平稳。

        综合路径规划与循迹仿真结果可以看出，Hybrid A*算法能够在有障碍物的仓储环境中生成满足差速运动学约束的可行路径，PID循迹控制器能够驱动差速底盘沿铺设引导线的路径稳定行驶。全局规划与局部循迹的两层架构验证了所设计方案在仿真层面的可行性：规划层负责路径决策，循迹层负责跟踪执行。需要指出的是，本仿真假定引导线已沿规划路径铺设完成，实际部署时需先将Hybrid A*的规划结果转化为地面引导标识。
    }
    }
\chapter{智能仓储移动机器人运动控制算法分析}
    \section{PID控制算法原理}
    \xfbody{PID控制器由比例（Proportional）、积分（Integral）和微分（Derivative）三个环节组成，是工业控制领域应用最为广泛的经典闭环控制方法，至今已有近百年的应用历史。PID控制器的基本思想是根据系统输出与期望值之间的偏差来构造控制量：比例环节对误差进行放大，提供快速的响应能力，但单独使用时会存在稳态误差；积分环节对误差进行累积，能够消除稳态误差，但可能导致响应变慢和超调增大；微分环节对误差的变化趋势进行预测，具有提前调节的作用，能够抑制超调与振荡，但对噪声较为敏感。三个环节相互配合，通过合理整定各自的增益系数，可以使系统同时具备较快的响应速度、较小的稳态误差和良好的动态品质。

在智能仓储移动机器人的运动控制中，PID控制器主要用于调节电机转速和机器人姿态响应。通过分别对比例增益$K_p$、积分增益$K_i$和微分增益$K_d$三个参数进行整定，可以改善系统对目标信号的跟踪速度和对干扰信号的抑制能力，使智能仓储移动机器人在不同工况下保持稳定的运动性能。

PID控制器的输出信号$u(t)$由三部分组成：

\begin{equation}
    u(t) = K_p e(t) + K_i \int_0^t e(\tau) d\tau + K_d \frac{de(t)}{dt}
    \label{eq:pid_control}
    \end{equation}
    其中，$e(t)$表示目标值与实际值之间的误差，$K_p$、$K_i$和$K_d$分别为比例增益、积分增益和微分增益。

    }
    \section{双轮差速PID控制器设计}
    \xfbody{

    针对双轮差速AGV的运动控制问题，本文采用双环PID级联控制架构。外环为位姿环，依据位置偏差和航向偏差计算期望线速度$v_{\mathrm{cmd}}$和期望角速度$\omega_{\mathrm{cmd}}$；内环为轮速环，由左右两个独立的PI控制器构成，负责跟踪外环下发的轮速指令并输出电机控制电压。内环响应频率约为外环的3$\sim$5倍，保证内环优先收敛后再由外环进行位姿修正。

    外环（位姿环）的输入为期望位姿与当前位姿之间的偏差，输出为线速度指令和角速度指令。以定点镇定为例，设目标位姿为$(x_t, y_t, \theta_t)$，当前位姿为$(x, y, \theta)$。当机器人距离目标较远时，期望航向指向目标点方向$\arctan(y_t-y, x_t-x)$，驱动机器人朝向目标点运动；当距离缩小至阈值以内后，期望航向切换为目标角度$\theta_t$，完成姿态对准。航向偏差经PID运算后输出角速度指令$\omega_{\mathrm{cmd}}$，距离偏差经比例运算后输出线速度指令$v_{\mathrm{cmd}}$。为避免大幅度转向时出现失控，当航向偏差超过$\pi/3$时将线速度缩减至正常值的$20\%$，优先完成转向。

    外环输出的$(v_{\mathrm{cmd}}, \omega_{\mathrm{cmd}})$经差速逆运动学转换为左右轮角速度指令：

    \begin{equation}
    \omega_{L,\mathrm{ref}} = \frac{v_{\mathrm{cmd}}}{r} - \frac{\omega_{\mathrm{cmd}}L}{2r},\quad
    \omega_{R,\mathrm{ref}} = \frac{v_{\mathrm{cmd}}}{r} + \frac{\omega_{\mathrm{cmd}}L}{2r}
    \label{eq:inverse_kinematics_ctrl}
    \end{equation}

    内环（轮速环）对左右轮各配置一个PI控制器。以左轮为例，控制律为：

    \begin{equation}
    u_L = K_p^v e_{\omega L} + K_i^v \int e_{\omega L}\,\mathrm{d}t,\quad e_{\omega L} = \omega_{L,\mathrm{ref}} - \omega_L
    \label{eq:wheel_pi}
    \end{equation}
    其中$K_p^v$和$K_i^v$分别为轮速环的比例增益和积分增益，$e_{\omega L}$为左轮速偏差，$u_L$为左电机电压指令。速度环采用PI而非PID的原因在于：微分项对编码器测量噪声较为敏感，在低速工况下容易引入电压抖动，恶化控制品质。

    双环PID控制器的各项参数见表\ref{tab:pid_params}。
    \begin{table}[H]
    \wuhao
    \centering
    \caption{双环PID控制器参数}
    \label{tab:pid_params}
    \begin{tabular}{cccc}
    \toprule
    控制器 & $K_p$ & $K_i$ & $K_d$ \\
    \midrule
    位姿环（航向PID） & 3.0 & 0.2 & 0.4 \\
    轮速环（左/右PI） & 2.0 & 12.0 & -- \\
    \bottomrule
    \end{tabular}
    \end{table}
    }
    \section{基本运动控制实现}
    \xfbody{

    控制器在STM32F103C8T6上以定时器中断方式周期性执行，中断周期$T_s=5\,\mathrm{ms}$。每个中断周期内的控制流程为：编码器脉冲计数→轮速计算→外环位姿PID→逆运动学解算→内环轮速PI→PWM输出。

    数字PID采用位置式算法，积分项做如下离散化处理：
    
    \begin{equation}
    u(k) = K_p e(k) + K_i T_s \sum_{j=0}^{k} e(j) + \frac{K_d}{T_s}\bigl[e(k) - e(k-1)\bigr]
    \label{eq:discrete_pid}
    \end{equation}

    为抑制积分饱和，轮速PI引入积分分离与遇限削弱积分两项措施。积分分离策略为：当轮速偏差$|e| > 8.0\,\mathrm{rad/s}$时清零积分项，仅保留比例控制以加速响应；当偏差落入阈值以内时重新激活积分，消除稳态误差。遇限削弱积分策略为：当PI输出超出电压限幅（$\pm12\,\mathrm{V}$）时停止正向积分累加，退回本次累加量，防止积分项在饱和状态下无意义增长。

    轮速测量利用SMC80S电机自带的$2500\,\mathrm{P/R}$增量式编码器，经STM32定时器编码器接口以四倍频模式采集，每转输出$10000$个脉冲。轮速估算采用M法（单位时间脉冲计数法）：$\omega = 2\pi \Delta N / (10000 \cdot T_s)$，其中$\Delta N$为$5\,\mathrm{ms}$内的脉冲增量。
    }
    \section{控制效果仿真验证}
    \xfbody{

    为验证双环PID控制器的有效性，在MATLAB环境下搭建差速AGV闭环仿真模型。仿真参数与实物设计保持一致：轮半径$r=0.075\,\mathrm{m}$，轮距$L=0.52\,\mathrm{m}$，电机机电时间常数$\tau=0.06\,\mathrm{s}$，电机增益$K_m=2.8\,\mathrm{(rad/s)/V}$，仿真步长$0.005\,\mathrm{s}$。设置四项测试场景：阶跃响应（验证速度跟踪能力）、定点镇定（验证位姿收敛能力）、直线跟踪（验证直行保持能力）和圆形轨迹跟踪（验证连续曲线跟踪能力）。

    \textbf{场景1——阶跃响应：}在$2.5\,\mathrm{s}$时施加幅度为$0.3\,\mathrm{m/s}$的线速度阶跃指令，同时保持角速度指令为零。仿真结果如图\ref{fig:pid_step_response}所示。速度响应无超调、无明显振荡，调节时间约$0.8\,\mathrm{s}$；稳态速度误差约为零，表明轮速PI能够快速且准确地跟踪参考指令。左右轮速指令在阶跃阶段保持一致，角速度始终维持在零附近，验证了直线行驶时两轮速度的同步性。电机电压在阶跃瞬间出现短暂峰值后迅速回落至稳态值，未触及饱和限。航向角全程保持$0^\circ$，未出现因轮速不一致导致的航向偏移。

    \begin{figure}[H]
    \centering
    \includegraphics[width=0.85\textwidth]{pid_step_response.png}
    \caption{阶跃响应仿真结果}
    \label{fig:pid_step_response}
    \end{figure}

    \textbf{场景2——定点镇定：}设定目标位姿为$(2.0, 1.0, \pi/4)$，即距起点约$2.24\,\mathrm{m}$、目标航向$45^\circ$。仿真结果如图\ref{fig:pid_point_stabilization}所示。机器人轨迹平滑收敛至目标点，终点位置误差为$(0.022, -0.007)\,\mathrm{m}$，终点航向误差为$-0.85^\circ$。运动过程分为两个明显的阶段：第一阶段机器人以前进转向方式朝向目标点运动，线速度随距离减小而平滑降低；第二阶段距离缩小至$8\,\mathrm{cm}$以下后，机器人主要进行原地姿态微调，直至航向误差收敛至$1^\circ$以内。

    \begin{figure}[H]
    \centering
    \includegraphics[width=0.85\textwidth]{pid_point_stabilization.png}
    \caption{定点镇定仿真结果}
    \label{fig:pid_point_stabilization}
    \end{figure}

    \textbf{场景3——直线跟踪：}设定目标直线为$y=0.6\,\mathrm{m}$，机器人以$0.25\,\mathrm{m/s}$的标称速度沿$x$轴正方向行驶，初始横向存在偏差。仿真结果如图\ref{fig:pid_line_tracking}所示，控制器根据横向偏差和航向偏差实时调整角速度指令，驱动机器人逐渐收敛至目标直线；约$5\,\mathrm{s}$后横向误差基本稳定在零附近，后续沿直线保持稳定行驶。稳态角速度RMSE为$0.0023\,\mathrm{rad/s}$，左右轮速在稳态下保持一致，航向角始终维持在$0^\circ$附近，表明控制器具备良好的直线保持能力。

    \begin{figure}[H]
    \centering
    \includegraphics[width=0.85\textwidth]{pid_line_tracking.png}
    \caption{直线跟踪仿真结果}
    \label{fig:pid_line_tracking}
    \end{figure}

    \textbf{场景4——圆形轨迹跟踪：}跟踪半径为$0.8\,\mathrm{m}$、周期为$16\,\mathrm{s}$的圆形轨迹，线速度约$0.31\,\mathrm{m/s}$。仿真结果如图\ref{fig:pid_circle_tracking}所示。实际轨迹与参考圆基本重合，轨迹圆度RMSE为$0.021\,\mathrm{m}$；线速度RMSE为$0.0088\,\mathrm{m/s}$，角速度RMSE为$0.0030\,\mathrm{rad/s}$。航向角呈线性增长，与参考航向一致，说明机器人以恒定角速度平稳完成圆周运动。左右轮速保持恒定差值，对应差速转向所需的稳态轮速分配。

    \begin{figure}[H]
    \centering
    \includegraphics[width=0.85\textwidth]{pid_circle_tracking.png}
    \caption{圆形轨迹跟踪仿真结果}
    \label{fig:pid_circle_tracking}
    \end{figure}

    四项场景的定量结果汇总如表\ref{tab:pid_sim_summary}所示。仿真结果表明，所设计的双环PID控制器在四种典型工况下均能实现较好的控制效果：定点镇定终点位置误差在厘米级、航向误差在$1^\circ$以内；连续轨迹跟踪的线速度和角速度RMSE均在$0.01\,\mathrm{m/s}$和$0.005\,\mathrm{rad/s}$量级，满足仓储AGV对基本运动控制精度的需求。

    \begin{table}[H]
    \wuhao
    \centering
    \caption{双环PID控制仿真结果汇总}
    \label{tab:pid_sim_summary}
    \begin{tabular}{ccc}
    \toprule
    测试场景 & 评价指标 & 结果 \\
    \midrule
    阶跃响应 & 速度稳态误差 & $\approx 0\,\mathrm{m/s}$ \\
    \multirow{2}{*}{定点镇定} & 终点位置误差 & $(0.022, -0.007)\,\mathrm{m}$ \\
    & 终点航向误差 & $-0.85^\circ$ \\
    直线跟踪 & 角速度RMSE & $0.0023\,\mathrm{rad/s}$ \\
    \multirow{3}{*}{圆形跟踪} & 线速度RMSE & $0.0088\,\mathrm{m/s}$ \\
    & 角速度RMSE & $0.0030\,\mathrm{rad/s}$ \\
    & 轨迹圆度RMSE & $0.021\,\mathrm{m}$ \\
    \bottomrule
    \end{tabular}
    \end{table}

    需要指出的是，本文仿真在理想地面条件下完成，未引入轮地摩擦系数扰动、电机死区、传感器噪声和通信延迟等因素。实际应用中AGV在负载变化、地面不平整等工况下的控制性能尚需通过实物实验进一步检验。此外，PID参数在仿真前整定后保持不变，对工况变化的适应能力有限，后续可引入模糊PID或模型预测控制等策略以提升鲁棒性。
    }
\chapter{结论}
    \section{全文工作总结}
    \xfbody{
    本文围绕智能仓储场景下举升式AGV小车的设计需求，从总体方案、机械结构、电气系统、运动学建模、路径规划与循迹控制、运动控制算法等方面展开了较为系统的研究工作，主要成果如下：

经过单轮驱动、差速驱动和全方位驱动三种方案的对比分析，确定了以双轮差速驱动为核心的举升式AGV总体方案，明确了关键技术指标包括：额定载重25kg、举升行程0.2m、最大行驶速度30m/min等。

机械结构层面，完成了移动底盘、驱动轮及减振机构、旋转升降机构和罩壳等核心部件的结构设计，并选取Q235A碳素结构钢作为支撑梁材料。通过有限元分析检验了支撑梁在满载工况下的最大应力为$4.891\,\text{MPa}$，远低于材料屈服强度235MPa，最大位移变形仅为$1.538\times10^{-3}\,\text{mm}$，最大应变为$2.268\times10^{-5}$，均处于安全范围内，印证了结构设计的合理性。

电气系统层面，选用STM32F103C8T6作为主控芯片，完成了SMC80S伺服电机与ZJPX115行星减速器的选型匹配，设计了TCRT5000红外传感器和HC-SR04超声波传感器的接口电路，利用电机自带编码器实现航位推算，同时完成了电源、时钟、SWD调试和复位等最小系统电路设计，为AGV的运动控制和环境感知奠定了硬件基础。

运动学建模层面，建立了双轮差速底盘的运动学模型，在MATLAB中对直线、圆弧、S形和8字形四类典型轨迹进行了仿真验证，分析了不同工况下的轮速分配规律和曲率变化特征。路径规划层面，采用Hybrid A*算法在$10\,\text{m} \times 8\,\text{m}$的仓储栅格地图上规划了无碰撞路径，经3134次迭代搜索得到长度为$10.71\,\text{m}$的可行路径。循迹控制层面，模拟5路TCRT5000红外传感器阵列进行PID循迹仿真，RMS横向误差为$0.8\,\text{cm}$，最大横向误差为$1.6\,\text{cm}$，验证了所设计方案在仿真层面的可行性。

运动控制算法层面，设计了位姿外环与轮速内环相结合的双环PID控制器，分别开展了阶跃响应、定点镇定、直线跟踪和圆形轨迹跟踪四类工况的仿真验证。结果显示：阶跃响应无超调、稳态速度误差接近零；定点镇定终点位置误差为厘米级、航向误差在$1^\circ$以内；直线跟踪稳态角速度RMSE为$0.0023\,\mathrm{rad/s}$；圆形轨迹跟踪线速度RMSE为$0.0088\,\mathrm{m/s}$、轨迹圆度RMSE为$0.021\,\mathrm{m}$，验证了所设计控制器在典型运动场景下的有效性。

}
    \section{不足和改进方向}
    \xfbody{
    尽管本文完成了AGV的总体方案设计、机械结构设计、电气系统设计、运动学建模和控制仿真等工作，但受时间和实验条件制约，研究仍存在不足。

在仿真验证层面，本文的分析和验证主要基于MATLAB/Simulink软件平台开展，尚未完成实物样机的搭建和实验测试。实际运行中可能存在的电机死区、轮胎打滑、传感器噪声和通信延迟等因素在仿真中未能充分体现，仿真结果与实际性能之间可能存在偏差。

在对比评估层面，本文目前仅给出了Hybrid A*+PID方案的仿真结果，尚未在统一地图、统一速度约束和统一评价指标下，与A*、DWA或其他常见控制策略开展定量对比。后续应补充基线实验与消融分析，以进一步增强结论的说服力。

在模型精度层面，本文建立的是二维平面运动学模型，主要描述位姿与速度之间的几何关系，对动力学特性（如惯性力、摩擦力矩）、负载变化（如载物台升降过程中的质心偏移）以及多传感器融合等问题考虑不够充分，整体仍属于简化模型，难以完全反映AGV在复杂工况下的真实运动状态。

在控制算法层面，本文采用的是经典PID控制方法，其参数在仿真前整定后保持不变。面对动态障碍物实时避让、多车协同调度、路径在线重规划等更复杂的任务场景，传统PID的自适应能力和学习能力有限，难以应对高度不确定的环境变化。

在软件功能层面，编码器闭环测速、多传感器信息融合、任务调度逻辑、无线通信组网以及上位机监控界面等软件功能尚未深入展开，距离形成完整的AGV控制系统仍有较大差距。

针对上述不足，后续可从多个方向推进：完成样机装配与联调，在真实环境中开展实物测试；引入动力学建模和轮地接触模型，提高仿真精度；结合模糊PID或模型预测控制等方法，提升跟踪精度和抗干扰能力；完善感知与通信模块，增强AGV的定位导航和多车协同能力。

}
    \section{应用前景展望}
    \xfbody{
    智能制造和无人仓储技术的快速推进，正促使AGV在工业生产和物流配送中扮演越来越重要的角色。本文研究的差速驱动举升式AGV具有结构简洁、控制方便、转弯灵活、制造成本较低等优点，尤其适用于中小型仓储、车间物料配送和轻载搬运等应用场景。

在仓储物流领域，这类AGV能够承担货物转运、工位配送、自动上下料和货架补货等多种任务。配合超声波传感器实现障碍物检测、红外传感器实现地面循迹后，AGV可在货架密集的仓储环境中自主行驶，有效减轻人工搬运的劳动强度，提高物料流转效率。对于电商仓储、医药物流、食品冷链和电子装配等对搬运柔性化和作业洁净度要求较高的场景，举升式AGV同样具有较强的实用价值和推广潜力。

从技术发展趋势来看，未来AGV将朝着更加智能化和协同化的方向演进。在感知层面，激光雷达、深度相机和毫米波雷达等高性能传感器的持续成熟，将使AGV具备更加精确和鲁棒的环境感知能力；在决策层面，边缘计算和人工智能芯片的引入，将使AGV能够在端侧完成复杂的路径规划和行为决策，降低对中央调度系统的依赖；在协同层面，结合仓储管理系统（WMS）和任务调度系统，多台AGV可以实现协同作业和动态任务分配，从单体搬运设备逐步升级为智能工厂和智能物流系统的关键执行单元。

对本设计而言，若后续完成实物样机的开发和系统集成调试，可在此基础上逐步拓展到自动充电管理、多任务动态调度、远程状态监控和多车协同导航等更高层次的应用功能，使举升式AGV真正成为智能物流系统的核心移动平台，为仓储自动化和智能制造的发展提供有力支持。

}



